\section{The Language-Factor Share, end to end}
\label{ch:lfs}

\subsection{The question}

When the same sentence is written in many languages and run through a model,
how much of the model's internal representation at a given layer is organized
by \emph{which language} it is, and how much by \emph{what it means}? LFS
answers with one number in $[0,1]$ per layer: near 1, the vectors cluster by
language whatever the sentence says; near 0, they cluster by meaning whatever
the language. Across layers the numbers form a depth profile.

\subsection{The grid}

Take $N$ sentences that exist in all $L$ languages with the same meaning.
NTREX-128 gives the WMT19 news test set in 128 languages; the profiles use
$N = 300$ and $L = 128$. Every cell (language $l$, sentence $c$) is one
string, and the grid is \emph{balanced}: every language has every sentence.
Balance is what makes the decomposition exact.

\subsection{From strings to vectors}

Run each string through the model with hidden states enabled. At one layer,
average the token vectors over the sentence's real tokens (not padding) in
fp32. That gives one vector $\mathbf h_{lc} \in \R^D$ per cell, and an
$L\times N\times D$ tensor per layer. Layer 0 is the embedding output; layer
$k$ is the residual stream after block $k$ (Chapter~\ref{ch:llm}).

\subsection{Standardize each coordinate, jointly}

For each coordinate $d$, subtract its mean and divide by its standard
deviation, both taken over all $L\times N$ cells together:
\[
z_{lcd} = \frac{h_{lcd} - \overline{h}_{\cdot\cdot d}}{s_d}.
\]
Why: a few coordinates carry values hundreds of times larger than the rest
and would decide the answer alone. Why jointly: standardizing each language
separately would subtract each language's mean, which is the quantity being
measured. After standardization every coordinate has the same total
variance, so LFS becomes a weighted average of per-coordinate language shares
in which no coordinate can dominate by scale.

\subsection{Three averages and two sums of squares}

\begin{figure}[h]\centering\includegraphics[width=\linewidth]{fig_lfs_grid.png}
\caption{Row means isolate language (they average over meaning); column means isolate meaning (they average over language).}\label{fig:grid}\end{figure}

Compute the grand mean $\bar{\mathbf z}$, the language means $\bar{\mathbf
z}_l$ (average over the $N$ sentences of language $l$), and the sentence means
$\bar{\mathbf z}_c$ (average over the $L$ translations of sentence $c$).
Averaging over sentences within a language cancels meaning and leaves
language identity; averaging over languages within a sentence cancels
language and leaves meaning. Then
\[
\SSl = N\sum_{l}\|\bar{\mathbf z}_l - \bar{\mathbf z}\|^2,\qquad
\SSc = L\sum_{c}\|\bar{\mathbf z}_c - \bar{\mathbf z}\|^2,\qquad
\SSr = \sum_{l,c}\|\mathbf z_{lc} - \bar{\mathbf z}_l - \bar{\mathbf z}_c + \bar{\mathbf z}\|^2,
\]
\[
\LFS = \frac{\SSl}{\SSl + \SSc}.
\]
The multipliers $N$ and $L$ are bookkeeping: each language mean stands for
$N$ cells and each sentence mean for $L$, so with those weights the three
terms add up exactly to the total spread of the grid. The residual is left out
of the ratio because LFS asks how the two \emph{systematic} effects divide; a
noisy model must not look language-agnostic by being noisy. The residual is
reported next to the ratio.

\subsection{A worked example for the whiteboard}

One coordinate, two languages, three sentences:
\begin{center}
\begin{tabular}{lrrr|r}
 & $s_1$ & $s_2$ & $s_3$ & language mean \\ \hline
English & 1.0 & 3.0 & 5.0 & 3.0 \\
German  & 2.0 & 4.0 & 6.0 & 4.0 \\ \hline
sentence mean & 1.5 & 3.5 & 5.5 & grand mean 3.5
\end{tabular}
\end{center}
$\SSl = 3\,[(3.0-3.5)^2 + (4.0-3.5)^2] = 1.5$;\quad
$\SSc = 2\,[(1.5-3.5)^2 + 0 + (5.5-3.5)^2] = 16$;\quad
every residual is zero (check: $1.0 - 3.0 - 1.5 + 3.5 = 0$);\quad
$\LFS = 1.5/17.5 = 0.086$.

Meaning dominates: sentences differ by a lot, languages by a constant
offset of one. Multiply every number by ten: every sum of squares is
multiplied by 100 and the ratio does not move. That is the scale invariance
that later becomes the collapse blind spot. (The REML variant in
\texttt{reports/lfs\_math.pdf} gives 0.11 on the same table because it
subtracts a noise estimate and divides by degrees of freedom; same data,
different estimator. On 14 real models the two rank models identically.)

\begin{exercise}{LFS from scratch, matching the repository code}
\begin{lstlisting}
def lfs(X):                                   # X: (L, N, D), already z-scored per coordinate
    L, N, D = X.shape
    mu = X.mean((0, 1)); mu_l = X.mean(1); mu_c = X.mean(0)
    ss_lang = N * ((mu_l - mu)**2).sum(); ss_con = L * ((mu_c - mu)**2).sum()
    ss_res = ((X - mu_l[:, None] - mu_c[None] + mu)**2).sum()
    return ss_lang / (ss_lang + ss_con), ss_lang, ss_con, ss_res
def zscore(H):                                # H: (L, N, D) raw pooled vectors
    flat = H.reshape(-1, H.shape[-1]); return (H - flat.mean(0)) / (flat.std(0) + 1e-9)
# illustration: language offsets G, concept vectors C, noise E
L, N, D = 12, 300, 128
G = rng.normal(0, 1.0, (L, 1, D)); C = rng.normal(0, 1.5, (1, N, D)); E = rng.normal(0, 0.5, (L, N, D))
H = G + C + E
print(lfs(zscore(H))[0])                       # a share between 0 and 1
print(lfs(zscore(3.7 * H))[0])                 # identical: uniform scale cancels
print(lfs(zscore(C + E))[0], (L-1)/(L+N-2))    # no language effect: near the null
\end{lstlisting}
Compare your function line by line with \texttt{code/pilot\_metrics.py:lfs}
and \texttt{src/analysis/rmfs\_profile.py:factor\_decomposition}. They are
the same estimator.
\end{exercise}

\subsection{What the profile looks like}

\begin{figure}[h]\centering\includegraphics[width=\linewidth]{fig_lfs_profiles_17.png}
\caption{Seventeen models on the 128-language, 300-sentence grid. Left: LFS by relative depth, four models highlighted. Right: dip depth, first-layer value minus minimum.}\label{fig:profiles}\end{figure}

Every one of the 17 primary models and the 19 expansion models has an
interior minimum with both endpoints higher (Figure~\ref{fig:profiles}).
Read it as three stages. Layer 0 is high because the embedding layer knows
the token identity and tokens are language-specific. The interior minimum is
where the model represents the same meaning most similarly across languages;
its depth runs from 0.039 (BLOOM-1.7B) to 0.336 (Qwen3-8B). The rise toward
the output is the model preparing to predict the next token, which is in a
specific language.

Two facts to state before anyone else does. The minimum is above 0.5 for
every model on this grid: even at the most shared layer, language main
effects exceed meaning main effects. Say ``more meaning-weighted relative to
its own endpoints'', never ``concept-dominated''. And the minimum's location
varies: most models place it between 35 and 60 percent depth, but Mistral-7B
(and two continued-pretraining models in the expansion) put a sharp minimum
at layer 2 followed by a plateau.

\begin{figure}[h]\centering\includegraphics[width=0.8\linewidth]{fig_lfs_decomp.png}
\caption{The same decomposition at three layers of Qwen3-8B. The residual (gray) is never in the ratio but always reported.}\label{fig:decomp}\end{figure}

\subsection{Properties, each with a short proof}

\textbf{Bounded.} Both sums of squares are non-negative, so $\LFS\in[0,1]$.

\textbf{Exact decomposition.} With $a_l = \bar{\mathbf z}_l - \bar{\mathbf
z}$, $b_c = \bar{\mathbf z}_c - \bar{\mathbf z}$, $r_{lc}$ the residual,
$\mathbf z_{lc} - \bar{\mathbf z} = a_l + b_c + r_{lc}$. On a balanced grid
$\sum_l a_l = 0$, $\sum_c b_c = 0$, and residuals sum to zero along every
row and column; expanding the squared norm and summing over the grid, every
cross term contains one of those zero sums. Hence total $=\SSl+\SSc+\SSr$.

\textbf{Uniform-scale invariance.} Multiply every vector by $k$: every sum of
squares is multiplied by $k^2$; the ratio is unchanged. After standardization
LFS is also invariant to any positive rescaling of individual coordinates.
Consequence: LFS cannot see all representations shrinking together
(Figure~\ref{fig:collapse}).

\textbf{Rotation invariance.} A rotation of the coordinate basis preserves
every norm in the sums, so the global LFS is unchanged. Which coordinates
carry the language energy is not invariant; hence no ``language neuron''
claims.

\textbf{The finite-grid null.} With no structure and isotropic noise of
variance $\sigma^2$ in $D$ coordinates, $\mathbb E[\SSl] = (L-1)D\sigma^2$
(language means average $N$ cells; $L-1$ degrees of freedom) and $\mathbb
E[\SSc] = (N-1)D\sigma^2$. Their ratio is
\[
\LFS_{\text{null}} = \frac{L-1}{L+N-2},
\]
independent of $D$: 0.298 at $(128, 300)$, 0.078 at $(128, 1500)$
(Figure~\ref{fig:null}). Three consequences: absolute values are not
comparable across grids, pooling schemes, or granularities; differences at a
fixed grid are fine because the noise term is common; and the null is not a
floor, since strong meaning structure with weak language structure drives
LFS below it. Verified on synthetic noise and on pseudo-language splits of
real representations to three decimals.

\begin{figure}[h]\centering\includegraphics[width=0.7\linewidth]{fig_lfs_null.png}
\caption{The pure-noise value depends on the grid. Compare LFS values only at a fixed grid.}\label{fig:null}\end{figure}

\textbf{Near-independence of evaluation size.} The noise term shrinks as a
share of the numerator as $N$ grows, so LFS drifts by $-0.006$ to $-0.010$
from $N=100$ to $N=1000$; a retrieval score moves $-0.065$ to $-0.125$ on the
same draws. This is the practical argument for LFS over retrieval scores when
comparing across studies.

\textbf{Differentiable.} LFS is a smooth function of the vectors, so it can
be a loss. That made the Aim 2 test possible; the test showed it should not be
one.

\subsection{The boundary and the failure mode}

LFS measures the \emph{additive} language component: how far each language's
mean sits from the center. Rotations, warps, and language-by-meaning
interactions land in the residual. The misalignment decomposition measured how much
that matters: offset plus scale explain 63 to 97 percent of held-out
cross-language misalignment, rotation under one percent, with the instrument
recovering injected rotations exactly. Language identity nonetheless stays
linearly decodable at every layer (probe accuracy 0.875 to 0.886, chance
0.042): variance share and separability are different quantities.

\begin{figure}[h]\centering\includegraphics[width=0.7\linewidth]{fig_lfs_collapse_synthetic.png}
\caption{Synthetic collapse: shrink every vector and the ratio is flat while the raw concept component vanishes.}\label{fig:collapse}\end{figure}

\begin{figure}[h]\centering\includegraphics[width=\linewidth]{fig_lfs_collapse_trained.png}
\caption{The same thing in a trained model. The word-MSE condition lost 95 percent of raw concept variance; LFS moved by 0.05 and MEXA ranked it first.}\label{fig:collapse_trained}\end{figure}

Own this before anyone raises it: in the Aim 2 word-alignment condition the
representation norm fell from 451 to 42 and raw concept variance from 13{,}791
to 699 while LFS moved from 0.445 to 0.493 and the retrieval metric ranked the
collapsed model first (Figure~\ref{fig:collapse_trained}). Every
scale-invariant measure shares this blind spot, by theorem. The fix is to
report the ratio with its raw components ($\SSl$, $\SSc$, residual, norm), so
that collapse shows up where it lives.

\subsection{What LFS relates to, and what it does not}

\begin{figure}[h]\centering\includegraphics[width=0.66\linewidth]{fig_r1_model_level.png}
\caption{Nineteen models: sentence-component share ($1-$LFS-VC) at the dip against median content transfer. Spearman 0.77 at the model level; 0.51 [0.19, 0.71] on the full confounder-adjusted panel.}\label{fig:r1}\end{figure}

\textbf{Content transfer (positive).} Give a model an English sentence to
predict, preceded by the previous sentence of the same document in another
language, versus a wrong-document sentence in that language; the improvement
from the right context is the content benefit. Across 19 models and 13
families, models with a more meaning-weighted dip layer show more content
transfer after controlling for data exposure, family, script, and fertility
(Figure~\ref{fig:r1}).

\textbf{Benchmark accuracy (null).} Against pooled multiple-choice accuracy on 33
models, LFS has no association once exposure is removed (0.04, interval
crossing zero); MEXA reaches 0.24. Two independent benchmarks agree on only about a
third of their per-language residual variation, which caps what any intrinsic
metric can predict there.

\textbf{Training objective (negative).} Minimizing LFS deepens the dip
quickly but through a shortcut that damages alignment and capability; across
18 trained condition-seed points lower LFS went with worse tail alignment. LFS is a
measurement, not a target.

\subsection{The family}

LFS answers one question. The other questions a multilingual representation
raises are answered by accompanying measures on the same grid, each validated on
its own question and never averaged into a single score:

\begin{center}\small
\begin{tabular}{p{4.6cm}p{5.2cm}p{4.6cm}}
\toprule
Question & Measure & Exposes \\
\midrule
Language vs.\ meaning share, per layer & LFS profile; dip depth & structure; recipe differences \\
Did scale or concept variance collapse? & raw $\SSl$, $\SSc$, residual, norm vs.\ a frozen reference & collapse a ratio hides \\
Is the cross-language map simple? & nested maps M0--M5, held-out error reduction & additive vs.\ rotational vs.\ nonlinear \\
Is the shared factor English? & latent-factor vs.\ English donor $R^2$ & English-centricity \\
Do the hardest sentences align? & AaR@10, worst-decile margin & tail failures that means hide \\
\bottomrule
\end{tabular}
\end{center}

MEXA is the published reference measure, not a family member. The proposal's own list
of candidate metric types (raw distances; mapping complexity; ``anything else
that emerges'') maps onto this table.

\subsection{How to say it}

\say{\textbf{Ten seconds.} LFS measures, at each layer, how much of the
systematic variation in a model's hidden states follows the language of a
sentence rather than its meaning, on a grid of matched translations.}

\say{\textbf{Thirty seconds.} Take the same 300 sentences in 128 languages,
run them through the model, and average token states into one vector per
sentence per language. Standardize coordinates jointly. The spread of the
language means measures language; the spread of the sentence means measures
meaning; LFS is the first divided by the sum. Across 33 models it falls to an
interior minimum and recovers, and how far it falls varies nine-fold and
tracks training recipe more than size. It must be reported with its raw
components, because the ratio cannot see everything shrinking at once.}

\say{\textbf{Two minutes.} Add: the exact two-way decomposition and why the
residual is excluded; the null $(L-1)/(L+N-2)$ and fixed-grid comparison;
values above 0.5 read relatively; the additive boundary measured at under one
percent rotation; the multilingual latent factor beating English in 113 to
124 of 127 languages; content transfer on 19 models at 0.51 [0.19, 0.71];
pooled benchmarks null with a measured ceiling; the collapse case; the family of
measures and the rule against averaging them.}
